\documentclass[10pt,a4paper]{amsart}
\usepackage[margin=19mm]{geometry}
\usepackage{amsmath,amssymb,mathtools}
\usepackage[scr=rsfs,cal=euler]{mathalfa}
\usepackage{xfrac}
\usepackage[unicode,hidelinks,bookmarks=false]{hyperref}
\newcommand{\ie}{i.e.\ }
\newcommand{\cf}{cf.\ }
\newcommand{\resp}{respectively\ }
\newcommand{\Subsection}{\subsection}
\providecommand{\nobreakdash}{\nobreak\hbox{-}}
\setlength{\emergencystretch}{2em}
\multlinegap0pt
\usepackage{amssymb}
\newtheorem{theorem}{Theorem}
\title{Trilinear characterizations of the Fourier extension conjecture on the paraboloid in three dimensions}
\author{Cristian Rios, Eric Sawyer}
\date{Source: arXiv:2506.03992v4 (CC BY 4.0)}
\begin{document}
\maketitle

\noindent\emph{Source and licence.} Trilinear characterizations of the Fourier extension conjecture on the paraboloid in three dimensions. Version arXiv:2506.03992v4 (CC BY 4.0), distributed by arXiv under the Creative Commons Attribution 4.0 International licence, http://creativecommons.org/licenses/by/4.0/, as recorded in the arXiv licence metadata for this exact version and shown on the abstract page https://arxiv.org/abs/2506.03992v4. Full licence notice: \texttt{licenses/arxiv-2506.03992-CC-BY-4.0.txt}.\par\smallskip

\noindent\textit{Editorial scope.} Counted unit: the article's theorem Alpert (source0.tex lines 1560--1568). The article's complete original proof of it is delivered with it (source0.tex lines 1767--1774, one \texttt{proof} environment of 8 lines, which the article places away from the statement). No mathematics is written, repaired or re-derived: the statement and the proof are the article's own lines, in the article's own notation, and no retained line is substituted or re-flowed.  Retained uncounted as the source-local context the proof needs: lines 187--191; lines 1697--1702; lines 1708--1711; lines 1730--1738.  Nothing was omitted from any retained span, so the package carries no omission record.  No retained line contains a citation, so the wrapper carries no bibliography and no bibliography entry.  No retained line contains a figure environment, an image-inclusion command or drawing code, so no image is delivered and the package declares no asset dependency.  Editorial formatting only, in addition: an AMS \texttt{amsart} preamble that restates the article's own declarations and macros used by the retained lines, namely the environments \texttt{theorem} on separate counters, together with the standard packages those lines need.  The retained environments print in this document as Theorem 1 in order of appearance, whereas the article prints the same items with the numbering of its own full document.  The complete native source of the article as distributed in the arXiv e-print for this exact version is bundled unchanged as \texttt{sources/179308/source0.tex}.
\begin{equation}
\left\Vert \mathcal{E}f\right\Vert _{L^{q}\left(  \mathbb{R}^{3}\right)
}\lesssim\left\Vert f\right\Vert _{L^{q}\left(  U\right)  }%
,\ \ \ \ \ \text{for }q>3.\label{FEC}%
\end{equation}

\begin{equation}
a\rightarrow\left\langle \delta_{a},\mu^{1}\ast\mu^{2}\right\rangle \text{ is
a smooth function of }a\text{ which is \emph{adapted to scale }}\frac{1}{\nu
}\min\left\{  2^{-s_{1}},2^{-s_{2}}\right\}  =\frac{1}{\nu}2^{-s_{2}%
},\label{ift adapted}%
\end{equation}

\begin{equation}
\mu^{1}\ast\mu^{2}\left(  a\right)  =\left\langle \delta_{a},\mu^{1}\ast
\mu^{2}\right\rangle .\label{density}%
\end{equation}

\begin{equation}
\left\vert \bigtriangleup_{I;\kappa}^{\eta}f\right\vert =\left\vert
\widehat{f}\left(  I\right)  h_{I;\kappa}^{\eta}\right\vert \lesssim
\ell\left(  I\right)  ^{-\left(  \frac{1}{p}-\frac{1}{p^{\prime}}\right)
}\left\Vert f\right\Vert _{L^{p}}\ell\left(  I\right)  ^{-1}\mathbf{1}%
_{\left(  1+\eta\ell\left(  I\right)  \right)  I}=\ell\left(  I\right)
^{-\frac{2}{p}}\left\Vert f\right\Vert _{L^{p}}\mathbf{1}_{\left(  1+\eta
\ell\left(  I\right)  \right)  I},\ \ \ \ \ 1<p<\infty.\label{precrude}%
\end{equation}

\begin{theorem}
\label{main Alpert}Let $0<\delta<1$ and $\kappa>\frac{20}{\delta}$. The
Fourier extension conjecture (\ref{FEC}) for the paraboloid in $\mathbb{R}%
^{3}$ holds \emph{if and only if} for every $q>3$ there is $\nu>0$ depending
only on $q$, such that the disjoint \emph{smooth Alpert }trilinear inequality
$\mathcal{A}_{\mathop{\rm disj}\nu}^{\kappa,\delta}\left(  \otimes
_{3}L^{\infty}\rightarrow L^{\frac{q}{3}};\varepsilon\right)  $ holds for all
$\varepsilon>0$.
\end{theorem}

\begin{proof}
[Proof of Theorem \ref{main Alpert}]Let $0<\delta<1$ and $\kappa>\frac
{20}{\delta}$ as in the statement of Theorem \ref{main Alpert}. Without loss
of generality we assume that $s_{1}<s_{2}<s_{3}$ and $R=2^{r}$, and we will
repeatedly use the inequalities (\ref{ift adapted}), (\ref{density}) and
(\ref{precrude}) to prove Theorem \ref{main Alpert} in three cases, followed
by a wrapup.
\end{proof}

\end{document}
